\partstart{Thiết kế giải pháp và quy trình học máy}
\subsection{Kiến trúc và hợp đồng dữ liệu}
Giải pháp tách thành bốn lớp: đọc và xác thực dữ liệu; tạo đặc trưng và tiền
xử lý; lựa chọn mô hình bằng cross-validation; đánh giá và xuất bằng chứng.
Ranh giới giữa các lớp giúp xác định bước nào học tham số và bước nào chỉ
kiểm tra hợp đồng. Tên cột, kiểu số, nhãn nhị phân và ngày hợp lệ phải được
xác thực trước khi huấn luyện. Khi hợp đồng sai, chương trình dừng thay vì
âm thầm sửa nhãn hoặc bỏ dòng.

\begin{figure}[!htbp]\centering
\fbox{\begin{minipage}{.88\textwidth}\centering
CSV và tệp mô tả $\longrightarrow$ Kiểm tra schema, nhãn, nhóm\\[7pt]
$\downarrow$\\[5pt]
Lấy mẫu nguyên đơn $\longrightarrow$ Chia train/test theo Order Id\\[7pt]
$\downarrow$\\[5pt]
Train: tạo đặc trưng $\to$ tiền xử lý trong fold $\to$ CV và tìm kiếm\\[7pt]
$\downarrow$\\[5pt]
Khóa lựa chọn bằng CV $\longrightarrow$ fit lại trên toàn train\\[7pt]
$\downarrow$\\[5pt]
Test lịch sử $\longrightarrow$ metric dòng/đơn, lỗi nhóm, bảng và hình
\end{minipage}}
\caption{Quy trình thực nghiệm; tập test nằm ngoài bước lựa chọn mô hình.}\label{fig:pipeline}
\end{figure}
Trong Hình~\ref{fig:pipeline}, kiểm toán toàn nguồn chỉ xác nhận chất lượng
và hợp đồng. Các biểu đồ khám phá phục vụ mô hình được dựng trên train sau
split. Các phép median, mode, scaling và học danh mục đều phải nằm trong
pipeline được fit lại trong từng fold. Việc chỉ chia train/test mà vẫn học
tiền xử lý trên toàn nguồn chưa đáp ứng điều kiện độc lập này.

Hàm \code{validate} yêu cầu nhãn chứa đúng hai giá trị số 0 và 1, mã đơn
không thiếu, trường số không vô hạn và ngày tạo đơn phân tích được. Điều
kiện hai lớp cũng được kiểm tra sau split và trong từng fold. Điều này ngăn
việc tính metric trên một tập không còn đủ lớp mà người đọc không biết.
Các quyết định thay đổi schema cần được thực hiện trước một protocol mới.

\subsection{Lựa chọn đặc trưng và phòng tránh rò rỉ}
Luồng A giữ bốn trường số: số ngày giao kế hoạch, số lượng mặt hàng, giá sản
phẩm và tỷ lệ chiết khấu. Chín trường phân loại gồm loại giao dịch, phương
thức giao, phân khúc khách, nhóm hàng, phòng ban, thị trường, vùng, quốc gia
và bang của đơn. Ngày tạo đơn chỉ được dùng để trích xuất lịch ở luồng B.
Mã đơn phục vụ chia nhóm, không được đưa vào vector dự đoán.

\begin{table}[!htbp]\centering\small
\caption{Các quyết định phòng rò rỉ trong danh sách đặc trưng.}\label{tab:leakage}
\begin{tabularx}{\linewidth}{|p{5cm}|Y|}\hline
\theadcell{Thông tin} & \theadcell{Quyết định và lý do}\\\hline
Delivery Status & Loại: trạng thái kết quả có liên hệ trực tiếp với nhãn.\\
Days for shipping (real) & Loại: chỉ biết sau vận chuyển, vi phạm thời điểm dự đoán.\\
Order Status; shipping date & Loại: thông tin sau tạo đơn hoặc thời điểm không bảo đảm.\\
Order Id và định danh cá nhân & Không làm đầu vào; mã đơn chỉ dùng gom nhóm.\\
Lợi nhuận, trường không trong allowlist & Loại bảo thủ do thời điểm/ngữ nghĩa chưa được xác nhận hoặc dư thừa.\\
Ngày tạo đơn & B trích tháng/thứ/giờ; không sử dụng nguyên timestamp.\\\hline
\end{tabularx}

\end{table}
Bảng~\ref{tab:leakage} phân biệt việc loại trường có kết quả tương lai với
việc tránh ghi nhớ định danh. Không phải mọi cột bị loại đều đã được chứng
minh gây leakage; một số bị loại vì chưa đủ bằng chứng về tính sẵn có.
\code{feature_audit.csv} lưu cả lý do và mức chắc chắn để không biến lựa
chọn bảo thủ thành kết luận tuyệt đối về bản chất từng trường.

Ngoài rò rỉ từ cột, các dòng cùng đơn xuất hiện ở hai phía split có thể tạo
liên hệ rất mạnh giữa train và test. Vì vậy, thiết kế chống leakage phải
kiểm soát đồng thời nội dung đặc trưng, đơn vị nhóm và nơi fit tiền xử lý.
Pipeline hỗ trợ kiểm soát phép biến đổi, nhưng không tự sửa một cách chia
nhóm sai hoặc hợp đồng thời điểm sai (xem \code{src/data.py} và \code{src/pipeline.py}).

\subsection{Lấy mẫu và chia train/test theo đơn}
Protocol \code{group-subset-v1} tính số dòng từng đơn, xáo trộn thứ tự đơn
với seed 42 rồi lấy tiền tố có tổng số dòng không vượt ngân sách 6.000.
Các đơn được giữ nguyên, không cắt lấy vài dòng của đơn cuối. Kết quả thực
tế là 5.999 dòng thuộc 2.113 đơn. Đây là mẫu theo nhóm có ràng buộc ngân
sách, không phải 6.000 dòng lấy ngẫu nhiên độc lập.

Sau lấy mẫu, \code{GroupShuffleSplit} tách khoảng 20\% nhóm làm test.
Với $G_{tr}$ và $G_{te}$ là tập mã đơn hai phía, điều kiện bắt buộc là
\begin{equation}
G_{tr}\cap G_{te}=\varnothing.
\end{equation}
Kết quả là 1.690 đơn/4.823 dòng train và 423 đơn/1.176 dòng test. Tỷ lệ
20\% áp dụng ở cấp nhóm; số dòng phụ thuộc kích thước từng đơn nên không
bắt buộc chia đúng 80/20 theo dòng. Chương trình kiểm tra giao nhóm bằng
tập hợp, thay vì chỉ dựa vào tên bộ chia (xem \code{src/data.py}).

Những đơn mới trong test vẫn có thể chia sẻ khách hàng hoặc sản phẩm với
train. Vì thế, split này trả lời câu hỏi dự đoán đơn chưa thấy trong cùng
quần thể lịch sử, chưa trả lời câu hỏi khách hàng mới, sản phẩm mới hoặc
tương lai. Muốn kiểm tra những mục tiêu đó cần định nghĩa nhóm hoặc thứ tự
thời gian khác và sử dụng dữ liệu đánh giá độc lập.

\subsection{Cross-validation theo nhóm và cách lựa chọn}
Trong train, \code{StratifiedGroupKFold} với ba fold, shuffle và seed 42
tạo danh sách cặp chỉ số dùng chung cho mọi ứng viên. Nhóm đơn không bị
chia cắt và bộ chia cố gắng cân bằng tỷ lệ lớp. Ràng buộc nhóm có ưu tiên,
nên không yêu cầu tỷ lệ lớp giữa các fold phải bằng tuyệt đối. Mỗi fold
được kiểm tra có đủ hai lớp và giao nhóm bằng rỗng.

Với cấu hình $\theta$ và $K=3$, điểm lựa chọn là
\begin{equation}
\overline F_1(\theta)=\frac{1}{K}\sum_{j=1}^{K}
F_1\!\left(y_{V_j},\hat f_{\theta,T_j}(X_{V_j})\right),\qquad
\theta^*=\arg\max_\theta\overline F_1(\theta).
\end{equation}
$T_j,V_j$ lần lượt là phần học và validation của fold $j$.
$\hat f_{\theta,T_j}$ bao gồm cả các phép tiền xử lý fit trên $T_j$.
F1 được tính cho lớp dương 1. Trung bình F1 từng fold không nhất thiết bằng
F1 khi gộp mọi dự đoán validation, vì F1 là một tỷ số phi tuyến.

\code{GridSearchCV} sử dụng \code{refit='f1'}: sau khi tìm cấu hình có
trung bình F1 cao nhất, nó fit lại pipeline trên toàn train. Tên ứng viên
thắng được lưu vào manifest trước vòng lặp dự đoán test. Test của sáu cấu
hình được báo cáo để so sánh đã định trước, không dùng để đổi người thắng.
Các quan sát lỗi test về sau chỉ tạo giả thuyết cho nghiên cứu tiếp theo.

Độ lệch chuẩn CV được lưu để mô tả dao động giữa ba fold, không là khoảng
tin cậy 95\%. Các fold có dữ liệu huấn luyện chồng lấn và cùng được dùng
chọn tham số, nên không thể coi ba điểm là ba thí nghiệm hoàn toàn độc lập.
Điểm CV tốt nhất còn có thể lạc quan do chọn cực đại trên nhiều cấu hình.

\subsection{Tiền xử lý A và B}
Với cột số $j$, giá trị thiếu được thay bằng median tính trên tập đang fit.
Sau đó dùng biến đổi $z_{ij}=(x_{ij}-\mu_j)/s_j$, trong đó $\mu_j,s_j$ là
trung bình và độ lệch chuẩn của tập học tương ứng. Không tính lại các tham
số này trên validation hay test. Chuẩn hóa giúp những trường khác đơn vị
không chi phối khoảng cách chỉ vì độ lớn số học.

Với trường phân loại, mode được học trên tập fit; one-hot biến mỗi danh mục
thành một chỉ báo nhị phân. \code{handle_unknown='ignore'} cho phép xử lý
danh mục chưa thấy, với các chỉ báo của trường đó bằng 0. Cơ chế này giúp
duy trì chiều đầu vào nhưng không bảo đảm dự đoán tốt cho danh mục lạ.
Ma trận thưa giúp hạn chế lưu trữ các số 0 trong không gian nhiều danh mục.

Luồng B giữ toàn bộ đặc trưng A và thêm tháng, thứ trong tuần, giờ tạo đơn.
Thứ dùng quy ước \code{dayofweek} của pandas: thứ Hai là 0, Chủ nhật là 6.
Ba trường được đưa vào nhánh số và chuẩn hóa như các trường số khác. Đây
là biểu diễn tuyến tính của biến lịch; không có mã hóa sin/cos trong code
hiện tại. Khoảng cách giữa giờ 23 và 0 vì vậy chưa phản ánh tính tuần hoàn.

Hai luồng có chung median, mode, scaling và one-hot; khác biệt trọng tâm
là ba đặc trưng lịch. Cây cũng sử dụng đầu vào đã scale để giữ cùng cấu
trúc pipeline. Với phép đổi thang đơn điệu của một trường số, thứ tự quan
sát được giữ nguyên, nên không gán lợi ích của cây cho StandardScaler khi
chưa có thí nghiệm tách riêng. Tương tự, không tuyên bố đã đo lợi ích của
imputation khi các cột đầu vào nguồn không có missing.

\subsection{K-Nearest Neighbors và liên hệ cấu hình}
K-NN dự đoán từ nhãn của $k$ quan sát gần nhất trong không gian đã tiền xử
lý. Gọi $\mathcal N_k(x)$ là tập láng giềng của điểm $x$, phiếu cho lớp $c$
được biểu diễn bằng $\sum_{i\in\mathcal N_k(x)}w_i\mathbf1(y_i=c)$.
Chế độ uniform đặt $w_i=1$; distance ưu tiên láng giềng gần hơn. Quy tắc
này giải thích vì sao biểu diễn dữ liệu và thang đo quan trọng đối với
K-NN \cite{knn}.

Hai khoảng cách được thử trong repository là
\begin{align}
d_2(x,z)&=\sqrt{\sum_{j=1}^{p}(x_j-z_j)^2},\\
d_1(x,z)&=\sum_{j=1}^{p}|x_j-z_j|.
\end{align}
$p$ là số chiều sau mã hóa. Một trường phân loại nhiều danh mục tạo nhiều
chiều one-hot; khi đó độ gần trong không gian mã hóa không nhất thiết là
độ tương đồng vận hành mà con người mong muốn. Những đơn cùng vùng có thể
vẫn rất khác về khả năng giao, trong khi một khác biệt danh mục hiếm có
thể tác động đáng kể tới thứ tự láng giềng.

Không gian tìm kiếm có $k\in\{5,15\}$, uniform/distance và Euclidean/
Manhattan, tổng cộng tám cấu hình cho mỗi luồng. Giá trị $k$ nhỏ cho phép
ranh giới linh hoạt nhưng nhạy với cục bộ; tăng $k$ làm trung bình trên
nhiều quan sát hơn. Không khẳng định $k=15$ tối ưu ngoài phạm vi lưới nhỏ
này. Kết quả được chọn còn phụ thuộc bộ đặc trưng, split và dữ liệu.

Code đặt \code{algorithm='brute'} và \code{n_jobs=1}. Với $n$ điểm học,
$m$ điểm cần dự đoán và $p$ chiều, phép tìm kiếm trực tiếp có chi phí tỷ lệ
xấp xỉ $mnp$ cho tính khoảng cách. Đây là lý do ngân sách dòng được giới
hạn trong thí nghiệm; báo cáo không suy ra thời gian chạy toàn bộ nguồn
bằng phép nhân tuyến tính đơn giản vì còn phụ thuộc bộ nhớ và xử lý thưa.

\subsection{Cây quyết định, độ thuần và regularization}
Cây phân loại chia không gian bằng các điều kiện trên một đặc trưng ở mỗi
nút. Với tỷ lệ lớp $p_{tc}$ tại nút $t$, hai tiêu chí trong lưới là
\begin{equation}
I_G(t)=1-\sum_c p_{tc}^2,\qquad
I_H(t)=-\sum_c p_{tc}\log_2 p_{tc}.
\end{equation}
Quy ước $0\log_2 0=0$. Một phép chia được đánh giá bằng mức giảm độ không
thuần có trọng số theo số mẫu ở hai nút con. Lá đưa ra nhãn theo phân bố
lớp của các mẫu tới lá; cây tìm quy tắc phân tách, không trực tiếp tối ưu
F1 ở từng nút \cite{tree}.

Lưới thử độ sâu tối đa 4, 10 hoặc không giới hạn; số mẫu tối thiểu một lá
5 hoặc 20; criterion là gini hoặc entropy. Có 12 cấu hình cho mỗi luồng.
Giới hạn độ sâu và tăng kích thước lá giảm khả năng tạo nhánh rất nhỏ chỉ
khớp dữ liệu học. Cây baseline dùng mặc định, lá tối thiểu 1 và không chặn
độ sâu, nên là đối chứng cho tác động của việc kiểm soát độ phức tạp.

Độ quan trọng impurity tổng hợp đóng góp giảm độ không thuần của đặc trưng
qua các nút, sau chuẩn hóa tổng. Giá trị cao nghĩa là đặc trưng được cây
dùng mạnh trong cách chia đã học; không phải hệ số nhân quả và không nói
rằng thay đổi đặc trưng sẽ làm thay đổi kết quả giao. Một đặc trưng có
importance bằng 0 trong cây nông vẫn có thể hữu ích với mô hình khác.

\subsection{Baseline và thước đo đánh giá}
Hai mốc so sánh có vai trò riêng. Cây mặc định kiểm tra việc dùng mô hình
linh hoạt nhưng chưa lựa chọn tham số. Dummy most-frequent học lớp đa số
trên train, sau đó luôn dự đoán lớp đó. Trong lần chạy này lớp đa số là 1.
Dummy không sử dụng thông tin đặc trưng để phân biệt các đơn, vì thế một
mô hình phức tạp không vượt mốc này cần được xem xét thận trọng.

Đặt TP là số dòng trễ được cảnh báo đúng, TN là dòng âm được dự đoán đúng,
FP là dòng âm bị cảnh báo và FN là dòng trễ bị bỏ sót. Với $N$ dòng test:
\begin{align}
\mathrm{Accuracy}&=\frac{TP+TN}{N},&
\mathrm{Precision}&=\frac{TP}{TP+FP},\\
\mathrm{Recall}&=\frac{TP}{TP+FN},&
F_1&=\frac{2TP}{2TP+FP+FN}.
\end{align}
Precision diễn tả tỷ lệ đúng trong số cảnh báo; Recall diễn tả phần dương
được phát hiện. F1 dùng để chọn mô hình trong protocol, nhưng bỏ qua TN và
không trực tiếp gắn chi phí vận hành. Code dùng \code{zero_division=0}
để xử lý trường hợp mẫu số bằng 0 một cách xác định.

Ngoài metric nhị phân, ROC-AUC đánh giá khả năng xếp hạng trên nhiều ngưỡng;
cột \code{pr_auc} trong code thực tế dùng \code{average_precision_score}.
Average precision là tổng precision có trọng số theo mức tăng recall,
khác với việc tích phân hình thang đường PR. Báo cáo giữ tên cột artifact
nhưng diễn giải đúng phép tính để không nhầm hai đại lượng (hàm \code{metrics} trong \code{src/run.py}).
Xác suất mô hình chưa được hiệu chuẩn.

\subsection{Đánh giá theo đơn và giả định thực nghiệm}
Với đơn $g$ chứa $n_g$ dòng, xác suất tổng hợp được định nghĩa
\begin{equation}
\bar p_g=\frac{1}{n_g}\sum_{i:g_i=g}p_i,\qquad
\hat y_g=\mathbf1(\bar p_g\ge0{,}5).
\end{equation}
Nhãn của đơn lấy từ nhãn chung sau khi kiểm tra tính nhất quán. Mỗi đơn
có trọng số bằng nhau khi tính metric. Quy tắc trung bình này đã được lưu
trong artifact, không được lựa chọn lại để tối đa hóa điểm test.

Thiết kế giả định schema và nhãn ổn định giữa train/test. Seed cố định
hỗ trợ tái lập, không chứng minh độ ổn định. Thử nghiệm chưa có temporal
holdout, chi phí lỗi, khoảng tin cậy nhóm hay đánh giá tác động can thiệp.
